\section*{ЗАКЛЮЧЕНИЕ}
\addcontentsline{toc}{section}{ЗАКЛЮЧЕНИЕ}

В ходе выполнения выпускной квалификационной работы был разработан web-сервис для размещения объявлений об услугах. Разработанная программная система предназначена для публикации услуг исполнителями, поиска подходящих предложений заказчиками, организации записи на свободные временные слоты и сопровождения взаимодействия между пользователями после создания бронирования.

В системе реализованы основные роли пользователей: неавторизованный пользователь, заказчик и исполнитель. Неавторизованный пользователь может просматривать каталог услуг и карточки предложений. Заказчик получает возможность выполнять поиск и фильтрацию услуг, добавлять предложения в избранное, создавать бронирования, вести переписку по заявке и оставлять отзывы. Исполнитель может публиковать и редактировать услуги, настраивать расписание, управлять недоступными датами, обрабатывать входящие заявки, изменять их статусы и оценивать взаимодействие с заказчиком.

Особое внимание при разработке было уделено механизму записи на услуги. В системе реализовано формирование свободных временных слотов на основе расписания исполнителя, длительности услуги и списка недоступных дат. При создании бронирования выполняется проверка доступности выбранного слота, что позволяет предотвратить повторную запись на одно и то же время.

Основные результаты работы:

\begin{enumerate}
	\item Проведен анализ предметной области. Выявлены основные проблемы ручной организации записи, определены ключевые участники системы и обоснована необходимость разработки web-сервиса.
	\item Сформулированы требования к программной системе. Определены функциональные возможности для заказчика, исполнителя и неавторизованного пользователя, а также требования к интерфейсу, надежности, информационной безопасности и программной совместимости.
	\item Выполнено проектирование web-сервиса. Разработана архитектура программной системы, спроектирована структура базы данных, определены основные пользовательские сценарии и элементы интерфейса.
	\item Реализована программная система средствами web-технологий. Разработаны модули регистрации и авторизации, управления профилем, каталога услуг, управления услугами исполнителя, расписания, формирования слотов, бронирования, обработки заявок, переписки, избранного и отзывов.
	\item Проведено системное тестирование разработанного web-сервиса. Проверены основные пользовательские сценарии, включая регистрацию, авторизацию, редактирование профиля, создание услуги, настройку расписания, формирование свободных слотов, поиск и фильтрацию услуг, бронирование, изменение статусов заявок, переписку, работу с избранным и отзывами.
\end{enumerate}

В результате выполнения работы были решены задачи, поставленные во введении. Разработанная программная система соответствует функциональным требованиям, сформулированным в техническом задании, и обеспечивает выполнение основных сценариев взаимодействия между заказчиком и исполнителем.

Разработанный web-сервис может использоваться как основа для дальнейшего развития системы размещения объявлений об услугах. В дальнейшем программная система может быть расширена за счет добавления онлайн-оплаты, уведомлений пользователей, более гибкой системы поиска и административной панели для модерации данных.